% !TeX root = main.tex
\newpage
\section{Methodology}

The gPCD framework integrates broad-phase occupancy generation directly
into the graphics pipeline. Rather than treating rendering as a
downstream consumer of simulation state, occupancy construction is
performed during graphics-pipeline execution within the vertex stage.
The resulting occupancy information is then used to construct collision
candidates for subsequent narrow-phase evaluation.

\subsection{Framework Overview}

\begin{figure}[h!]
\centering
\begin{verbatim}

Particle State
       |
       v

+--------------------------------------+
|          GRAPHICS PIPELINE           |
|                                      |
|  Vertex Shader                       |
|      - Corner Generation             |
|      - Cell Determination            |
|      - Occupancy Construction        |
|                                      |
|  Fragment Shader                     |
|      - Rendering                     |
+--------------------------------------+
               |
               v
      Candidate Collision Pairs
               |
               v
+--------------------------------------+
|          COMPUTE PIPELINE            |
|                                      |
|  Narrow-Phase Evaluation             |
|                                      |
+--------------------------------------+
               |
               v
          Collision Set
\end{verbatim}
\caption{High-level data flow through the gPCD framework. Broad-phase occupancy
generation is integrated directly into graphics-pipeline execution,
where particle geometry is used to construct collision candidates. The
resulting candidate pairs are subsequently evaluated within the compute
pipeline during narrow-phase collision detection.}
\label{fig:gpcd_dataflow}
\end{figure}


Figure~\ref{fig:gpcd_dataflow} illustrates the high-level organization
of the gPCD framework. The framework consists of a graphics-pipeline
broad phase and a compute-shader narrow phase. Particle state is first 
processed by the graphics pipeline, where graphics-pipeline occupancy 
generation constructs cell occupancy information from particle geometry.
The resulting occupancy data are stored within a uniform cell array and
used to construct candidate collision pairs. These candidates are then
evaluated within a compute shader during narrow-phase collision
detection.
\begingroup
\begin{figure*}[ht]
\centering
\hspace{0.3in}
	\begin{subfigure}[b]{0.3\textwidth}
		\includegraphics[width=\textwidth]{CellSpansA.png}
		\subcaption[]{Corner Duplicates}
		\label{fig:CellSpansA}
	\end{subfigure}
\hspace{0.3in}
	\begin{subfigure}[b]{0.3\textwidth}
		\includegraphics[width=\textwidth]{CellSpansB.png}
		\subcaption[]{Collision Duplicates}
		\label{fig:CellSpansB}
	\end{subfigure}
\hspace{0.3in}
\caption[TITLE:Cell Spans]{Two-dimensional illustration of occupancy and collision duplication
arising from corner-based occupancy construction. (a) Multiple corners
may map to the same cell, resulting in duplicate occupancy
contributions. (b) A single collision may be represented within
multiple neighboring cells when interacting particles span a cell
boundary. The three-dimensional implementation follows identically
using the eight corners of the particle AABB.}
\Description[TITLE:Cell Spans]{Two-dimensional illustration of occupancy and collision duplication
arising from corner-based occupancy construction. (a) Multiple corners
may map to the same cell, resulting in duplicate occupancy
contributions. (b) A single collision may be represented within
multiple neighboring cells when interacting particles span a cell
boundary. The three-dimensional implementation follows identically
using the eight corners of the particle AABB.}
\label{fig:CellSpans}
\end{figure*}
\endgroup


\subsection{Graphics-Pipeline Occupancy Generation}

Broad-phase occupancy generation is implemented within the vertex stage
of the graphics pipeline. Each particle independently evaluates its
occupancy contribution and determines the cell locations required for
broad-phase construction. The resulting occupancy information is stored
within the cell-array structure and subsequently used to construct
candidate collision pairs for narrow-phase evaluation.


\subsection{Corner-Based Occupancy Construction}

Each particle is represented by an axis-aligned bounding box (AABB).
Occupancy generation is performed by evaluating the eight corners of the
AABB and determining the cell locations intersected by the particle.
The resulting corner locations define the occupancy contribution of the
particle within the cell array.

Let $(x,y,z)$ denote the floating-point particle position and let
$r$ denote the particle radius. The eight corner locations are
constructed as

\begin{equation}
\mathbf{x}_c =
\mathbf{x}
+
(\pm r,\pm r,\pm r),
\end{equation}

where $c \in \{0,\ldots,7\}$ identifies one of the eight AABB corners.

Each corner location is subsequently converted to a cell coordinate and
used to determine particle occupancy within the cell array.

The corner-based formulation guarantees that particle occupancy can be
determined through a fixed number of evaluations independent of particle
count. Because the particle radius is constrained such that a particle
may intersect at most eight neighboring cells, occupancy construction
remains bounded by the eight AABB corner locations.

\subsection{Cell Addressing and Flattening}

Each corner location is mapped to a cell coordinate within the uniform
cell array. Let $(x,y,z)$ denote a floating-point corner location and
let $(X,Y,Z)$ denote the corresponding unsigned integer cell
coordinates. Cell coordinates are obtained by rounding the corner
location to the nearest integer cell center,

\begin{equation}
X = \mathrm{round}(x),
\qquad
Y = \mathrm{round}(y),
\qquad
Z = \mathrm{round}(z).
\end{equation}

Cell coordinates are restricted to positive integer values such that
the valid cell domain begins at $(1,1,1)$ rather than $(0,0,0)$.
Coordinates containing a zero-valued component are invalid.
This convention prevents corner evaluations associated with particles
located adjacent to domain boundaries from producing negative cell
coordinates during occupancy construction.

The resulting cell coordinate triplet is converted into a flattened cell
index for storage within the occupancy structure. Let $S$ denote the
side length of the cubic cell array. The flattened index is

\begin{equation}
I = X + YS + ZS^2.
\end{equation}

The flattened index serves as the primary cell identifier throughout
the gPCD collision-detection pipeline. In addition to directly
addressing the corresponding occupancy list within the cell array, the
flattened representation reduces cell-reference storage from a
three-component coordinate tuple to a single integer identifier. Once
computed, the index is used throughout occupancy construction,
candidate generation, and narrow-phase evaluation without requiring
reconstruction of the original cell coordinates.

\subsection{Potentially Colliding Set Construction}

The populated cell array constitutes the potentially colliding set
(PCS) used during narrow-phase evaluation. Each occupied cell contains
the identifiers of particles whose corner locations map to that cell.
Consequently, occupancy generation directly constructs the PCS without
requiring a separate candidate-construction stage.

Because a particle may occupy multiple cells and a collision may be
represented within multiple neighboring occupancy lists, occupancy
duplication and collision duplication can arise within the PCS.
Figure~\ref{fig:occupancy_duplication} illustrates these effects. The
resulting PCS is subsequently traversed during narrow-phase collision
evaluation.


\subsection{Narrow-Phase Evaluation}

The candidate pairs generated during broad-phase construction are
evaluated within a compute shader during narrow-phase collision
detection. Each candidate pair is processed independently to determine
whether a collision exists between the corresponding particles.

Because occupancy construction may generate duplicate candidate pairs,
narrow-phase evaluation operates on the potentially colliding set
generated during candidate construction. The resulting collision set is
subsequently used by the particle dynamics system.
